\setcounter{table}{0}
\renewcommand{\thetable}{S\arabic{table}}

\section{Supplementary Reproducibility Information}

This supplement provides reproducibility details that are not fully expanded in the main text, including the exact task prompt, output inclusion rules, evaluated output counts, rating anchors, analysis artifacts, and reproducibility boundaries.

\section{S1. Unified Task Prompt}

All models were asked to complete the same biomedical research-analysis task:

\begin{verbatim}
Based on public transcriptomic data, construct a multi-gene signature
for predicting immunotherapy response in non-small cell lung cancer
(NSCLC), and explore the role of programmed cell death mechanisms,
including ferroptosis, cuproptosis, and pyroptosis, in immunotherapy
resistance.

Please output a research-analysis plan and complete data-analysis
workflow with a level of complexity comparable to a journal article
with an impact factor (IF) of approximately 5 in the corresponding year.

The content should cover public dataset selection, cohort design,
endpoint definition, data preprocessing, differential expression
analysis, candidate gene screening, model construction, validation
strategy, immune microenvironment analysis, mechanistic interpretation,
key figures and tables, and manual review points.
\end{verbatim}

For the skill-augmented strategy, the same task was executed in an artificial intelligence (AI) agent environment with autonomous access to the medical research skill package. The agent was allowed to select relevant skills without a fixed human-imposed sequence.

\section{S2. Output Inclusion Rules and Evaluated Output Counts}

Outputs were retained if they produced reviewable research-analysis material, including:

\begin{itemize}
\item complete research-analysis outputs;
\item partial research reports with sufficient task-relevant content;
\item stage summaries;
\item archived traces that contained sufficient task-relevant content for human review.
\end{itemize}

Outputs were not excluded solely because the final artifact format differed across model or agent runs, provided that the archived material contained enough content for blinded human review. This rule was used to preserve real-world agentic-output variability.

\begin{table}[htbp]
\centering
\caption{Evaluated outputs by model and generation strategy.}
\label{tab:supp-output-counts}
\small
\begin{tabularx}{\linewidth}{Y C{1.8cm} C{1.8cm} Y}
\toprule
Model & Native-AI outputs & Skill-augmented outputs & Notes \\
\midrule
GPT-5.4 & 2 & 2 & Included in model-specific descriptive analysis \\
Claude Sonnet 4.6 & 1 & 2 & Included in model-specific descriptive analysis \\
GLM-5.1 & 1 & 2 & Included in model-specific descriptive analysis \\
DeepSeek-V4 Pro & 2 & 2 & Included in model-specific descriptive analysis \\
Kimi K2.6 & 1 & 2 & Included in model-specific descriptive analysis \\
MiniMax-M2.7 & 2 & 2 & Included in model-specific descriptive analysis \\
\bottomrule
\end{tabularx}
\end{table}

The final analysis included 21 anonymized outputs in total: 9 native-AI outputs and 12 skill-augmented outputs. Anonymous output identifiers were used during blinded review, but the complete identifier-to-output archive is not included in the arXiv source package because the full generated materials require additional de-identification before public release.

\section{S3. Rating Scale Anchors and Reviewer Design}

All Likert-scale items used a 1-7 scale. Each anonymized output received two non-expert biomedical reviewer ratings and two blinded expert ratings. Reviewers did not receive model names, generation strategies, platform configuration, or the mapping between anonymous output IDs and experimental conditions.

\begin{table}[htbp]
\centering
\caption{Rating scale anchors.}
\label{tab:supp-rating-anchors}
\small
\begin{tabularx}{\linewidth}{Y Y Y Y}
\toprule
Construct & Anchor 1 & Anchor 4 & Anchor 7 \\
\midrule
Quality, clarity, completeness, feasibility, usability & Very poor / strongly disagree & Neutral or partially acceptable & Excellent / strongly agree \\
Workflow coherence & Incoherent or poorly integrated & Partially coherent & Highly coherent and well integrated \\
Expert methodological quality & Methodologically inappropriate & Partially appropriate with limitations & Methodologically strong and appropriate \\
Perceived risk & Very low perceived risk & Moderate or uncertain risk & Very high perceived risk \\
\bottomrule
\end{tabularx}
\end{table}

For quality and workflow constructs, higher scores indicate better evaluation. For perceived-risk constructs, higher scores indicate greater perceived methodological or interpretive risk.

\section{S4. Displayed Analysis Artifacts}

The following analysis artifacts are displayed in the compiled arXiv PDF. Aggregate CSV files used for analysis are retained in the source package for reproducibility but are not listed here as display artifacts because CSV files are not rendered as manuscript content during submission.

\begin{table}[htbp]
\centering
\caption{Displayed analysis artifacts.}
\label{tab:supp-analysis-artifacts}
\small
\begin{tabularx}{\linewidth}{Y Y}
\toprule
File & Purpose \\
\midrule
\path{figure0_quality_control_flow.png} & Evaluation dataset quality-control flow figure \\
\path{figure1_overall_quality_by_strategy.png} & Overall expert and non-expert quality figure \\
\path{figure2_secondary_outcomes_by_strategy.png} & Expert methodological quality and perceived risk figure \\
\path{figure3_model_specific_skill_effects.png} & Model-specific skill-minus-native effects figure \\
\bottomrule
\end{tabularx}
\end{table}

\section{S5. Reproducibility Boundaries}

The study is reproducible at the level of task prompt, model labels, strategy definitions, output inclusion criteria, anonymized output IDs, rating constructs, and analysis outputs. It is not claimed to be exactly reproducible at the level of proprietary model sampling internals, hidden provider-side model revisions, or OpenClaw deployment internals that were not exported in the available run archives.
